\section*{Введение}
\addcontentsline{toc}{section}{Введение}

Современные CI/CD-процессы к настоящему моменту достаточно надёжно
контролируют \emph{функциональную} корректность программного обеспечения:
модульные и интеграционные тесты, статический анализ и линтеры являются
стандартной, повсеместно принятой практикой и в большинстве промышленных
пайплайнов блокируют слияние кода при обнаружении регрессии
функциональности. \emph{Производительность}, однако, в подавляющем
большинстве пайплайнов остаётся вне поля автоматической проверки:
регрессия времени выполнения, объёма потребляемой памяти или числа
аллокаций типично обнаруживается постфактум — либо вручную при
эпизодическом профилировании, либо только в продакшене, когда деградация
уже сказалась на пользователях. Исторически интерес к профилированию
производительности возникал прежде всего как реакция на инциденты в
эксплуатации — от ранних систем трассировки вызовов до промышленных
инфраструктур непрерывного профилирования дата-центров (Google-Wide
Profiling~\cite{ren2010}) и их современных открытых аналогов (Grafana
Pyroscope~\cite{pyroscope-doc}, Parca~\cite{parca-doc}) — то есть решал
задачу мониторинга уже выложенной системы, а не предотвращения попадания
регрессии в неё. Актуальность настоящей работы состоит в переносе этой
идеи на более ранний этап жизненного цикла — момент pull request, до
слияния кода с основной веткой, — где стоимость исправления регрессии на
порядки ниже, чем после развёртывания в продакшене.

Целью работы является разработка и экспериментальная апробация прототипа
системы, интегрирующей профилирование производительности Go-приложений в
CI/CD-пайплайн для автоматического обнаружения регрессий производительности
на раннем этапе. Задачи, решение которых обеспечивает достижение этой
цели, — анализ существующих инструментов и подходов, формализация
критерия обнаружения регрессии, проектирование и программная реализация
архитектуры прототипа, а также экспериментальная апробация на управляемых
сценариях — подробно сформулированы по результатам аналитической части в
разделе~\ref{sec:razdel1}.

Теоретическая значимость работы заключается в формализации критерия
обнаружения регрессии производительности, явно разделяющего
статистическую значимость отклонения (критерий Манна --- Уитни или
Уэлча~\cite{welch1947}) и его практическую величину, а также в
разграничении понятия "непрерывное профилирование" на процессный смысл
(встроенное в каждый прогон CI) и эксплуатационный смысл (постоянный
мониторинг production-системы), что позволяет корректно позиционировать
данную работу относительно существующих промышленных решений
(раздел~\ref{sec:1.2}). Практическая значимость — в том, что предложенный
прототип реализован как переносимый набор скриптов и CI-пайплайн,
применимый к произвольному Go-проекту с бенчмарками без изменения кода
слоя сравнения, что подтверждено архитектурным решением о разделении
приложения и слоя сравнения (раздел~\ref{sec:razdel3}).

Работа выполнена в рамках учебно-исследовательской работы студента и не
является частью более широкого кафедрального или производственного
проекта.

Пояснительная записка состоит из пяти основных разделов. Раздел~1
посвящён анализу инструментов профилирования и бенчмаркинга в экосистеме
Go и сравнительному анализу существующих подходов к интеграции
профилирования в CI/CD, по результатам которого формулируются цель и
задачи работы. В разделе~2 формализуются модель процесса профилирования в
CI/CD-пайплайне и статистический критерий обнаружения регрессии,
сочетающий практическую и статистическую значимость отклонения. Раздел~3
описывает проектирование архитектуры прототипа: обоснование выбора модели
жизненного цикла и инструментальных средств, состав и структуру основных
компонентов системы. В разделе~4 излагается программная реализация
прототипа — тестовое Go-приложение, реализация критерия сравнения с
эталоном, настройка CI/CD-пайплайна и подходы к тестированию. Раздел~5
представляет результаты экспериментальной апробации прототипа на четырёх
управляемых сценариях регрессии, включая оценку корректности обнаружения
и накладных расходов профилирования. В заключении подводятся итоги работы
и намечаются направления дальнейших исследований.
